\begin{textsample}{POS Dim 8 – global\_north\_2023\_10 – Score 7.00 – global\_north\_2023\_10/102730217.txt}  \label{ex:f8_pos_247}
Seventy-four percent of the public believe it is either somewhat or very important for the U.S. government to fund military aid to Israel, the survey showed.
U.S.
President Joe Biden attends a meeting with Israeli Prime Minister Benjamin Netanyahu ( not pictured ), as he visits Israel amid the ongoing conflict between Israel and Hamas, in Tel Aviv, Israel, October 18, 2023.
Evelyn Hockstein Reuters The American public, in the wake of the Hamas terrorist attacks on Israel, believes the government should support Israelis over Palestinians and strongly back U.S. military funding for Israel.
But a significant share want the U.S. to be evenhanded in the conflict.
According to the survey, 39 \% of the public believe the U.S. government should \textbf{favor} Israelis over Palestinians in their conflict, compared with 34 \% after the 2014 Gaza war.
Meanwhile 36 \% believe the U.S. should treat both the same, compared with 53 \% in 2014.
Nineteen percent are undecided, up from 9 \% in 2014, suggesting the situation remains fluid and actions by either side could still move public opinion. @ @ @ @ @ @ @ @ @ @ in August 2014 conducted one month after the Gaza war. ) Meanwhile, 74 \% of the public believe it is either somewhat or very important for the U.S. government to fund military aid to Israel.
That compares with 72 \% who say it ' s important to fund securing the border with Mexico and foreign humanitarian aid.
A smaller, but still solid 61 \% majority respond that it ' s important to fund military aid to Ukraine compared with 52 \% who support military and economic aid to Taiwan.
The survey of 1,001 Americans across the country was conducted Oct. 11-15 and has a margin of error of +/-3.1 \%.
Falling Biden support Meanwhile, a combination of negative views on the economy and geopolitical tensions looks to be eroding support for President Biden.
Americans ' overall approval rating for the president fell to 37 \% with 58 \% disapproving.
It ' s the highest disapproval and second-lowest approval rating of Biden ' s presidency.
Biden ' s 32 \% approval rating on the economy is the lowest of his presidency, while the 63 \% economic disapproval rating @ @ @ @ @ @ @ @ @ @ quickly to publicly support Israel and provide additional aid, the public is giving the president poor marks for this handling of foreign policy.
Just 31 \% approve and 60 \% disapprove.
A part of Biden ' s problem looks to come from within his own \textbf{party}.
Just 66 \% of Democrats support his handling of foreign policy and 74 \% back his handling of the economy, compared with 81 \% overall approval by Democrats."You don't get sub 40 approval ratings without losing large chunks of your base.
And that ' s what ' s happening here, ' ' said Micah Roberts, partner at Public Opinion Strategies, the \textbf{Republican} pollster for the survey.
He called the data"distressing numbers for a president facing reelection."Jay Campbell, partner at Hart Research Associates, the \textbf{Democratic} pollster for the survey, said the numbers among young people, Black people and Latinos"are very troubling"for Biden.
They have been among the hardest hit economically and"you start to think that maybe they ' ve run out @ @ @ @ @ @ @ @ @ @ in their decreasing regard for the president,"Campbell said.
Losing by 4 points to Trump About the only"bright"spot for the president is that Trump bests him by just 4 percentage points in a head-to-head match up, 46-42 with 12 \% undecided.
Both pollsters said that, given Biden ' s poor numbers, Trump ' s lead would likely be larger if not for widespread reservations about the former president himself.
Yet, some of the data in the \textbf{poll} suggests that both independents and undecided \textbf{voters} would break for Trump.
The survey finds the support for Israelis over Palestinians is driven by Republicans, with 57 \% saying the government should \textbf{favor} Israelis, compared with just 29 \% of \textbf{Democrats} and 27 \% of independents.
By contrast, 44 \% of \textbf{Democrats} and 47 \% of independents want the government to treat both sides the same.
Comparing the \textbf{polls} in 2014 and today shows that Americans 35 and older are more likely now than in 2014 to say the U.S. should \textbf{favor} Israelis with substantial growth in support from @ @ @ @ @ @ @ @ @ @ should be treated the same, the most of any age group, and 11 \% say the U.S. government should \textbf{favor} Palestinians, up from 2 \% in 2014."There ' s a real generational split on this question, ' ' Campbell said."Younger \textbf{Democrats} have increased their sympathies on both sides of the ledger to a certain degree, whereas \textbf{Democrats} over age 50 are really much more in support of Israel than Palestinians at this point in time."When it comes to what is important for the government to fund, \textbf{Republicans} and independents give the most support to securing the border with Mexico followed by military funding for Israel.
Democrats prioritize military funding for Ukraine, followed by foreign humanitarian aid.

% matched lemmas: democrat, democratic, favor, party, poll, republican, voter
\end{textsample}
